% SOURCE_AUTHORITY: sga5_fr_workpass.tex, lines 11350--11401 (LF-normalized unit).
% SOURCE_SHA256: 791F4EFFC5E02832D5D77ED03518C8156D6F07E4C8238B03545DB93D883FBB28
\subsection*{9.8}

Demonstremos agora 9.6. Para isso, vejamos primeiro a igualdade
\[
v_*f^*i_*(y)=v_*j_*(g^*(y)c_{d-1}(\check F)),
\tag{9.8.1}
\]
e depois demonstraremos que $v_*$ é injetiva. De (9.7) (iii) resulta
\[
v_*f^*i_*(y)=f_1^*u_*i_*(y)-(j_1)_*[g_1^*i^*i_*(y)],
\]
Isto é, tendo em conta (9.2),
\[
v_*f^*i_*(y)=f_1^*(i_1)_*(y)-(j_1)_*g_1^*(y\,c_d(\check N)).
\tag{9.8.2}
\]
Por outro lado, tem-se imediatamente
\[
v_*j_*(g^*(y)c_{d-1}(\check F))
=(j_1)_*\bigl[k_*k^*g_1^*(y)c_{d-1}(\check F)\bigr],
\]
Ou seja, pela fórmula de projeção para $k$,
\[
v_*j_*(g^*(y)c_{d-1}(\check F))
=(j_1)_*\bigl[g_1^*(y)k_*(c_{d-1}(\check F))\bigr].
\]
Utilizando (9.7)(i), deduz-se
\[
v_*j_*(g^*(y)c_{d-1}(\check F))
=(j_1)_*(g_1^*(y)c_d(\check E))-(j_1)_*g_1^*(y\,c_d(\check N)).
\tag{9.8.3}
\]
Comparando (9.8.2) e (9.8.3), vemos que (9.8.1) equivale a
\[
f_1^*(i_1)_*(y)=(j_1)_*(g_1^*(y)c_d(\check E)),
\tag{9.8.4}
\]
isto é, à igualdade correspondente a (9.6.2) quando se substitui $X$ por $X\times\mathbb P^1$. Por (9.1.2), tem-se
\[
(j_1)_*(g_1^*(y)c_d(\check E))=(j_1)_*s_*(y)=\beta_*\delta_*(y).
\]
Mas é imediato (cf. $(D_1)$) que $\delta_*=f_2^*\gamma_*$, de onde:
\[
(j_1)_*(g_1^*(y)\,c_d(\check E))=\beta_*f_2^*\gamma_*(y).
\tag{9.8.5}
\]
Por outro lado, tem-se evidentemente
\[
f_1^*(i_1)_*(y)=f_1^*\alpha_*\gamma_*(y),
\tag{9.8.6}
\]
de modo que (9.8.4) resulta da comparação de (9.8.5), (9.8.6) e (9.7)(ii).
